\GalSourcePageBoundary{073}{L0072}{44}{44}

permutations de l'un des groupes partiels, une fonction qui n'est
invariable par aucune substitution.)

Soit $\theta$ cette fonction des racines.

Opérons sur la fonction $\theta$ une des substitutions du groupe total
qui ne lui sont pas communes avec les groupes partiels. Soit $\theta_{1}$ le
résultat. Opérons sur la fonction $\theta_{1}$ la même substitution, et
soit $\theta_{2}$ le résultat, et ainsi de suite.

Comme $p$ est un nombre premier, cette suite ne pourra s'arrêter
qu'au terme $\theta_{p-1}$; ensuite l'on aura
\[
\GalControlReading{\theta_{p}=\theta_{0},\qquad \theta_{p+1}=\theta_{1}}{W09-CD002},
\]
et ainsi de suite.

Cela posé, il est clair que la fonction
\GalSourceErrorMark{W09-SE004}
\[
\bigl(\theta+\alpha\theta_{1}
        +\alpha^{2}\theta_{2}+\ldots
        +\alpha^{p-1}\theta_{p-1}\bigr)^{p}
\]
sera invariable par toutes les permutations du groupe total et, par
conséquent, sera actuellement connue.

Si l'on extrait la racine $p^{\text{ième}}$ de cette fonction, et qu'on l'adjoigne
à l'équation, alors, par la proposition IV, le groupe de
l'équation ne contiendra plus d'autres substitutions que celles des
groupes partiels.

Ainsi, pour que le groupe d'une équation puisse s'abaisser par
une simple extraction de racine, la condition ci-dessus est nécessaire
et suffisante.

Adjoignons à l'équation le radical en question; nous pourrons
raisonner maintenant sur le nouveau groupe comme sur le précédent,
et il faudra qu'il se décompose lui-même de la manière
indiquée, et ainsi de suite, jusqu'à un certain groupe qui ne contiendra
plus qu'une seule permutation.

\emph{Scolie.} --- Il est aisé d'observer cette marche dans la résolution
connue des équations générales du quatrième degré. En effet, ces
équations se résolvent au moyen d'une équation du troisième
degré, qui exige elle-même l'extraction d'une racine carrée. Dans
la suite naturelle des idées, c'est donc par cette racine carrée
qu'il faut commencer. Or, en adjoignant à l'équation du quatrième
degré cette racine carrée, le groupe de l'équation, qui contenait
en tout vingt-quatre substitutions, se décompose en deux qui n'en
contiennent que douze. En désignant par $a,b,c,d$ les racines,
